% ============================================================
\section{Pipeline Pseudocode}
\label{app:pipeline}
% ============================================================

The following pseudocode summarises the full execution sequence of
the \texttt{crns-sitechar} pipeline (cf.\ \texttt{main.py}).
Each numbered step corresponds to a function call that may read from
or write to the on-disk cache.

\begin{lstlisting}[language=Python, caption={CRNS site characterization
  pipeline — top-level pseudocode (main.py).},
  label={lst:pipeline}]
# -------------------------------------------------------
# Input parameters (from config or CLI)
# -------------------------------------------------------
LAT, LON          = 46.4, 11.3        # decimal degrees
THETA_V_INIT      = 0.20              # m^3/m^3, design moisture
N_SL              = 2800              # cph, detector reference
OUTPUT_DIR        = "output/"

# -------------------------------------------------------
# Step 1-2  Terrain
# -------------------------------------------------------
dem, dx, dy = download_dem(LAT, LON, buffer=3*r86_approx())
slope, aspect, horizon = compute_terrain_derivatives(dem, dx, dy)
twi = compute_twi(dem, slope, cache_dir=OUTPUT_DIR)   # cached

# -------------------------------------------------------
# Step 3   Preliminary footprint parameters
# -------------------------------------------------------
r86_prelim = r86(THETA_V_INIT, P=site_pressure(ELEV))
z86_prelim = z86(THETA_V_INIT, rho_b=1.2, lw=0.0)

# -------------------------------------------------------
# Step 4   Land-cover and OSM
# -------------------------------------------------------
worldcover = download_worldcover(LAT, LON, r86_prelim)
osm        = download_osm(LAT, LON, r86_prelim)

# -------------------------------------------------------
# Step 5   Soil properties (SoilGrids)
# -------------------------------------------------------
soil = download_soilgrids(LAT, LON)
rho_b   = soil['bulk_density_gcc']
clay_gg = soil['clay_gg']
soc_gkg = soil['soc_gkg']
lw      = 0.097 * clay_gg + 0.033     # Kohli 2021

# -------------------------------------------------------
# Step 6   Updated footprint with actual lw and rho_b
# -------------------------------------------------------
z86_act = z86(THETA_V_INIT, rho_b, lw)
r86_act = r86(THETA_V_INIT, P=site_pressure(ELEV))

# -------------------------------------------------------
# Step 7   kappa_topo  (3-D ray-casting)
# -------------------------------------------------------
kappa_topo = compute_kappa_topo(dem, z86_act, rho_b,
                                n_theta=18, n_phi=72)

# -------------------------------------------------------
# Step 8   kappa_muon  (horizon integral)
# -------------------------------------------------------
kappa_muon = compute_kappa_muon(horizon)

# -------------------------------------------------------
# Step 9   kappa_lulc  (W(r)-weighted f_H)
# -------------------------------------------------------
lulc_result = get_lulc(worldcover, osm, dem, dx, dy,
                       r86_act, theta_v_init=THETA_V_INIT)
kappa_lulc  = lulc_result['wc_kappa']

# -------------------------------------------------------
# Step 10  Site fluxes and N0
# -------------------------------------------------------
kappa_tot = kappa_topo * kappa_muon * kappa_lulc
flux = compute_site_fluxes(LAT, LON, ELEV, N_SL,
                           lw=lw, soc_gkg=soc_gkg,
                           rho_b=rho_b,
                           kappa_tot=kappa_tot)
N0  = flux['N0']

# -------------------------------------------------------
# Step 11  Atmospheric / snow seasonality  (ERA5)
# -------------------------------------------------------
crns_corr = get_crns_corrections(LAT, LON,
                                 lai=vegetation['lai_mean'])

# -------------------------------------------------------
# Step 12  Geological characterisation  (Macrostrat)
# -------------------------------------------------------
geology = get_geology(LAT, LON)

# -------------------------------------------------------
# Step 13  Sampling plan
# -------------------------------------------------------
sampling = compute_sampling_plan(r86_act, P=site_pressure(ELEV))
plot_sampling_plan(sampling, dem, dx, dy, dist_grid)

# -------------------------------------------------------
# Step 14  Assemble and write report
# -------------------------------------------------------
results = {
    'kappa_topo': kappa_topo,
    'kappa_muon': kappa_muon,
    'kappa_lulc': kappa_lulc,
    'kappa_tot':  kappa_tot,
    'N0':         N0,
    'lw':         lw,
    'soc_gkg':    soc_gkg,
    'crns_corrections': crns_corr,
    'geology':    geology,
    'sampling':   sampling,
    ...
}
write_report(results, OUTPUT_DIR)
save_json(results, OUTPUT_DIR / "results.json")
\end{lstlisting}

Figure~\ref{fig:flowchart} illustrates the data flow between modules.

\begin{figure}[H]
  \centering
  % The flowchart is generated at run time by the pipeline itself
  % and saved to OUTPUT_DIR/pipeline_flowchart.pdf.
  \includegraphics[width=0.95\linewidth]{pipeline_flowchart}
  \caption{Data-flow diagram of the CRNS site characterization
    pipeline.
    Rectangles denote processing steps; cylinders denote on-disk
    caches; rounded rectangles denote external data sources.
    Arrows indicate data dependencies.}
  \label{fig:flowchart}
\end{figure}

\paragraph{Cache strategy.}
Each expensive step writes its result to a uniquely named file in
\texttt{OUTPUT\_DIR} (e.g.\
\texttt{dem\_mosaic.npz}, \texttt{soilgrids\_0\_30cm.json},
\texttt{twi\_cache\_<hash>.npz}).
At the start of each step the pipeline checks whether the output
file already exists and, if so, loads it directly.
This means that a run interrupted halfway can be resumed from the
last successful step with no data re-download.

\paragraph{Configuration.}
All tunable parameters (site coordinates, $N_\mathrm{sl}$,
$\thetav^\mathrm{ref}$, ray-casting resolution, ring fractions,
download time-out) are collected in a single block at the top of
\texttt{main.py} and are documented in Table~\ref{tab:variables}.
No external configuration file is required for standard use.
